% Part: first-order-logic
% Chapter: syntax-and-semantics
% Section: substitution

\documentclass[../../../include/open-logic-section]{subfiles}

\begin{document}

\olfileid{fol}{syn}{ext}

\olsection{Extensionaliteit}

\begin{explain}
Extensionaliteit, soms ook relevantie genoemd, kan informeel als
volgt worden uitgedrukt: de enige factoren die van invloed zijn op de
vervulling van !!{formula}~$!A$ in !!a{structure}~$\Struct M$ onder
!!a{variable} bedeling~$s$, zijn de omvang van het !!{domain} en de
waarden die~$\Struct M$ en~$s$ toekennen aan de taalelementen die
daadwerkelijk in~$!A$ voorkomen.

Een onmiddellijk gevolg van extensionaliteit is het volgende. Als twee
!!{structure}s~$\Struct M$ en~$\Struct M'$ overeenstemmen op alle
taalelementen die in een zin~$!A$ voorkomen en hetzelfde domein hebben,
dan moeten~$\Struct M$ en~$\Struct M'$ ook overeenstemmen over de vraag
of~$!A$ zelf waar is.
\end{explain}

\begin{prop}[Extensionaliteit]
  \ollabel{prop:extensionality}
  Zij $!A$ !!a{formula}, en zijn $\Struct M_1$ en $\Struct M_2$
  !!{structure}s met $\Domain{M_1} = \Domain{M_2}$, en $s$ een
  bedeling op $\Domain{M_1} = \Domain{M_2}$. Als
  $\Assign{c}{M_1} = \Assign{c}{M_2}$, $\Assign{R}{M_1}=\Assign{R}{M_2}$,
  en $\Assign{f}{M_1} = \Assign{f}{M_2}$ voor elke !!{constant}~$c$,
  elk relatiesymbool~$R$ en elke !!{function}~$f$ die in~$!A$
  voorkomen, dan geldt $\Sat{M_1}{!A}[s]$ dan en slechts dan als
  $\Sat{M_2}{!A}[s]$.
\end{prop}

\begin{proof}
  Bewijs eerst met inductie naar~$t$ dat voor elke term die in de
  formule voorkomt, $\Value{t}{M_1}[s] = \Value{t}{M_2}[s]$. Bewijs vervolgens
  de propositie met inductie naar~$!A$ en gebruik de zojuist bewezen
  bewering in het basisgeval van de inductie (waarin~$!A$ atomair is).
\end{proof}

\begin{prob}
Werk het bewijs van \olref[fol][syn][ext]{prop:extensionality}
in detail uit.
\end{prob}

\begin{cor}[Extensionaliteit voor !!^{sentence}s]
  \ollabel{cor:extensionality-sent}
  Zij $!A$ !!a{sentence} en zijn $\Struct{M_1}$ en $\Struct{M_2}$ als in
  \olref{prop:extensionality}. Dan geldt $\Sat{M_1}{!A}$ dan en slechts
  dan als $\Sat{M_2}{!A}$.
\end{cor}

\begin{proof}
Dit volgt uit \olref{prop:extensionality} met
\olref[ass]{cor:sat-sentence}.
\end{proof}

Bovendien hangen de waarde van een term en de vraag of !!a{structure}
!!a{formula} vervult uitsluitend af van de waarden van de deeltermen.

\begin{prop}\ollabel{prop:ext-terms}
Zij $\Struct M$ !!a{structure}, zijn $t$ en $t'$ termen, en zij $s$ een
bedeling. Dan geldt $\Value{\Subst{t}{t'}{x}}{M}[s] =
\Value{t}{M}[\Subst{s}{\Value{t'}{M}[s]}{x}]$.
\end{prop}

\begin{proof}
Met inductie naar~$t$.
\begin{enumerate}
\item Als $t$ een constante is, zeg $t\ident c$, dan is
  $\Subst{t}{t'}{x} = c$, en $\Value{c}{M}[s] = \Assign{c}{M} =
  \Value{c}{M}[\Subst{s}{\Value{t'}{M}[s]}{x}]$.

\item Als $t$ een andere variabele dan~$x$ is, zeg $t \ident y$, dan
  is $\Subst{t}{t'}{x} = y$, en geldt
  $\Value{y}{M}[s] =
  \Value{y}{M}[\Subst{s}{\Value{t'}{M}[s]}{x}]$, omdat
  $\varAssign{s}{\Subst{s}{\Value{t'}{M}[s]}{x}}{x}$.

\item Als $t \ident x$, dan is $\Subst{t}{t'}{x} = t'$. Maar
  $\Value{x}{M}[\Subst{s}{\Value{t'}{M}[s]}{x}] = \Value{t'}{M}[s]$
  volgens de definitie van~$\Subst{s}{\Value{t'}{M}[s]}{x}$.

\item Als $t \ident \Atom{f}{t_1,\dots,t_n}$, dan geldt:
\begin{multline*}
  \Value{\Subst{t}{t'}{x}}{M}[s] \\
  \begin{aligned}[b]
& = \Value{\Atom{f}{\Subst{t_1}{t'}{x}, \dots, \Subst{t_n}{t'}{x}}}{M}[s]\\
& \qquad    \text{ volgens de definitie van $\Subst{t}{t'}{x}$}\\
& = \Assign{f}{M}(\Value{\Subst{t_1}{t'}{x}}{M}[s], \dots,
    \Value{\Subst{t_n}{t'}{x}}{M}[s])\\
    & \qquad  \text{ volgens de definitie van $\Value{\Atom{f}{\dots}}{M}[s]$}\\
& = \Assign{f}{M}(\Value{t_1}{M}[\Subst{s}{\Value{t'}{M}[s]}{x}], \dots,
   \Value{t_n}{M}[\Subst{s}{\Value{t'}{M}[s]}{x}])\\
& \qquad    \text{ volgens de inductiehypothese}\\
& = \Value{t}{M}[\Subst{s}{\Value{t'}{M}[s]}{x}]
    \text{ volgens de definitie van $\Value{\Atom{f}{\dots}}{M}[\Subst{s}{\Value{t'}{M}[s]}{x}]$}
  \end{aligned}
\end{multline*}
\end{enumerate}
\end{proof}

\begin{prop}\ollabel{prop:ext-formulas} Zij $\Struct M$
!!a{structure}, $!A$ !!a{formula}, $t'$ een term, en $s$ een bedeling.
Dan geldt $\Sat{M}{\Subst{!A}{t'}{x}}[s]$ dan en slechts dan als
$\Sat{M}{!A}[\Subst{s}{\Value{t'}{M}[s]}{x}]$.
\end{prop}

\begin{proof}
Opgave.
\end{proof}

\begin{prob}
Bewijs \olref[fol][syn][ext]{prop:ext-formulas}.
\end{prob}

\begin{explain}
  De strekking van
  \cref{fol:syn:ext:prop:ext-terms,fol:syn:ext:prop:ext-formulas} is
  als volgt. Stel dat we een term~$t$ of !!a{formula}~$!A$ en een
  term~$t'$ hebben, en dat we de waarde van $\Subst{t}{t'}{x}$ willen
  kennen of willen weten of $\Subst{!A}{t'}{x}$ wordt vervuld in
  !!a{structure}~$\Struct M$ onder !!a{variable} bedeling~$s$.
  Dan kunnen we eerst de substitutie uitvoeren en vervolgens de
  waarde of vervulling onder $\Struct{M}$ en~$s$ beschouwen. We kunnen
  ook eerst de waarde~$m = \Value{t'}{M}[s]$ van~$t'$ in
  $\Struct{M}$ onder~$s$ bepalen, de !!{variable} bedeling veranderen
  in~$\Subst{s}{m}{x}$, en vervolgens de waarde van~$t$ in
  $\Struct{M}$ onder~$\Subst{s}{m}{x}$ beschouwen, of nagaan of
  $\Sat{M}{!A}[\Subst{s}{m}{x}]$.
  \Cref{fol:syn:ext:prop:ext-terms,fol:syn:ext:prop:ext-formulas}
  garanderen dat het antwoord op beide manieren hetzelfde is.
\end{explain}

\end{document}
